%------------------------------------------------------------------------------%
\section{Integração por partes}
\label{sec:Integracao_por_partes}

%------------------------------------------------------------------------------%
\begin{proposicao}{Integração por Partes}{pro:integracao-partes}
  Dadas funções integráveis $p$ e $q$, podemos escrever a seguinte relação
  entre as integrais indefinidas
  \[
  \int u\,dv = uv - \int v\,du \,.
  \]
  De modo equivalente temos a relação entre as integrais definidas
  \[
  \int_{-\infty}^\infty p\,dq = pq \, \evalat{-\infty}{\infty} - \int_{-\infty}^\infty q\,dp \,.
  \]
\end{proposicao}
%------------------------------------------------------------------------------%

\hrulefill


\textbf{Quando usar?} Normalmente, usamos integração por partes quando queremos
integrar o produto de duas funções onde uma delas é a derivada de alguma função
conhecida, isto é
\[
  \int f(x) g'(x) dx
\]

\textbf{Exemplos:}
\begin{itemize}
  \item Em $\int x \sin(x) dx$, temos um produto de funções onde $\sin(x)$, por
      exemplo, é a derivada de alguém que conheço, isto é, de $-cos(x)$.

  \item $\int x lnx dx$, aqui também temos um produto de funções e $x$ é a
      derivada de $\frac{x^2}{2}$

  \item $\int (x^2+1) e^{-x} dx$, aqui temos um produto de duas funções e
      $e^{-x}$ é a derivada de $-e^{-x}$.

  \item Nossa intenção ao usarmos integração por partes é obtermos uma integral
  mais simples que aquela de partida.
\end{itemize}

%------------------------------------------------------------------------------%
\begin{teorema}{Fórmula de Integração por partes}{teo:formula-integral-partes}
  Nem sempre é imediato o cálculo de integrais.
  Por isso, assim como em derivadas,
  precisamos de algumas técnicas para o cálculo dessas
  \begin{equation}
    \label{eq:formula_partes-1}
    \int f(x) g'(x) dx= f(x) g(x)- \int g(x) f'(x)dx
  \end{equation}
  Chamando $u=f(x)$ e $v=g(x)$, então $du=f'(x)dx$ e $dv=g'(x)dx$
  \begin{equation}
    \label{eq:formula_partes-2}
    \int u dv=uv- \int v du
  \end{equation}
\end{teorema}
%------------------------------------------------------------------------------%

Para demonstrar a fórmula de integração por partes,
lembremos da regra do produto para derivadas
\[
  (f(x)g(x))'= f'(x)g(x)+f(x)g'(x)
\]
Integrando ambos os lados da equação com relação a $x$, temos
\[
  \int (f(x)g(x))' dx= \int f'(x)g(x) dx+ \int f(x)g'(x) dx
\]
ou seja,
\[
  (f(x)g(x)) = \int f'(x)g(x) dx+ \int f(x)g'(x) dx
\]
Isolando o termo que queremos temos que
\[
  \int f(x)g'(x) dx= f(x)g(x)-\int f'(x)g(x) dx+ \int f(x)g'(x) dx
\]

\begin{itemize}
  \item É útil chamar $u=f(x)$ e $dv=g'(x)dx$ e usar o padrão
    \begin{align*}
      u  & = f(x)    &  dv & = g'(x)dx \\
      du & = f'(x)dx &  v  & = g(x)
    \end{align*}
    quando temos uma integral do tipo $\int f(x)g'(x)dx$,
    pois facilita a memorização da fórmula de integração por partes.

  \item Na fórmula de integração por partes, $g$ é qualquer
    antiderivada de $g'$, pegamos a com constante $0$ por simplicidade.

  \item Você não deve escolher $u$ e $dv$ de qualquer maneira.
    Essa escolha deve ser feita de modo a obter uma integral mais
    simples do que a que tínhamos inicialmente. Com a prática de
    exercícios essa escolha fica simples.
\end{itemize}

%------------------------------------------------------------------------------%
\begin{example} \label{ex:ex1partes}
  Encontre a Integral \(\int x \sin(x)dx\)

  \solution
  \emph{Solução 1} -- usando a fórmula \eqref{eq:formula_partes-1}:
  Vamos escolher $f(x)=x$ e $g'(x)= \sin(x)$.
  Logo $f'(x)=1$ e $g(x)= -\cos(x)$ (Lembre-se que para $g$
  podemos escolher qualquer antiderivada de $g'$.).
  Assim, usando a fórmula \ref{f1}, obtemos
  \begin{align*}
    \int x\sin(x)dx
    & = f(x)g(x)- \int g(x)f'(x)dx \\
    & = x(-\cos(x))- \int (-\cos(x))dx \\
    & = -x \cos(x) + \int \cos(x) dx \\
    & = -x \cos(x) + \sin(x)  + C
  \end{align*}

  \emph{Solução 2} -- usando a fórmula \eqref{eq:formula_partes-2}:
  Usando o padrão
  \begin{align*}
    u  & =  x & dv & = \sin(x) dx \\
    du & = dx & v  & = -\cos(x)
  \end{align*}
  Assim, fórmula \eqref{eq:formula_partes-2} nos dá
  \begin{align*}
    \int x\sin(x)dx
    & = uv- \int vdu \\
    & = x\big(-\cos(x)\big)- \int \big(-\cos(x)\big)dx \\
    & = -x \cos(x) + \int \cos(x) dx \\
    & = -x \cos(x) + \sin(x)  + C
  \end{align*}
\end{example}
%------------------------------------------------------------------------------%

Você poderia perguntar: Eu não poderia escolher $u= \sin(x)$ e $dv=x $, uma vez
que conseguimos encontrar uma antiderivada de modo fácil para $x$? Como
comentamos anteriormente, a escolha de $u$ e $dv$ deve ser feita de modo a
obter uma integral mais simples. Caso tivéssemos feito essa escolha chegaríamos
na integral
\[
    \int x \sin(x) dx= \sin(x) \frac{x^2}{2}- \frac{1}{2} \int x^2\cos(x)dx
\]
que é uma integral mais difícil que aquela que começamos.

%------------------------------------------------------------------------------%
\begin{example}
  Calcule \(\int \ln(x) dx\)
  \solution
  Primeiro note que podemos rescrever
  \[
    \int \ln(x) dx= \int \ln(x) \cdot 1 dx
  \]
  Usando o padrão temos
  \begin{align*}
    u  & = \ln(x)        & dv & = 1 dx \\
    du & = \frac{1}{x}dx & v  & = x
  \end{align*}
  Integrando por partes chegamos que
  \begin{align*}
    \int \ln(x) dx
    & = x \ln(x)- \int  x \frac{dx}{x} \\
    & = x \ln(x)-\int dx \\
    & = x \ln(x)-x +C
  \end{align*}
\end{example}
%------------------------------------------------------------------------------%

%------------------------------------------------------------------------------%
\begin{example}
  Encontre \(\int t^2 e^t dt\)
  \solution
  Note que $t^2$ torna-se mais simples quando derivamos enquanto que
  $e^t$ permanece inalterada quando integramos ou derivamos.
  Assim, escolhemos
  \begin{align*}
    u  & = t^2  & dv & = e^t dt \\
    du & = 2tdt & v  & = e^t
  \end{align*}
  Logo, a fórmula \eqref{eq:formula_partes-2} de integração por partes
  resulta em
  \begin{equation}
    \label{eq:exemplo_r}
    \int t^2 e^t dt= t^2 e^t-2 \int te^t dt
  \end{equation}
  Observe que a integral $ \int te^t dt $ não é imediata.
  Nesse caso, temos que usar novamente Integração por partes.
  Escolhemos,
  \begin{align*}
    u  &= t  & dv & = e^t dt \\
    du &= dt & v  & = e^t
  \end{align*}
  Assim,
  \[
    \int te^t dt
    = t e^t - \int e^t dt
    = t e^t - e^t + C
  \]
  Substituindo esse resultado na equação \ref{eq:exemplo_r}, obtemos
  \begin{align*}
    \int t^2 e^t dt
    & = t^2 e^t - 2\int te^tdt  \\
    & = t^2 e^t - 2(te^t-e^t+C) \\
    & = t^2 e^t - 2te^t+2e^t+C_1
  \end{align*}
  onde $C_1=2C$.
\end{example}
%------------------------------------------------------------------------------%

%------------------------------------------------------------------------------%
\begin{example}
  Calcule \(\int e^x \sin(x) dx\)
  \solution
  Usando integração por partes
  \begin{align*}
    u  & = e^x   & dv & = \sin(x) dx \\
    du & = e^xdx & v  & = -\cos(x)
  \end{align*}
  Assim
  \begin{equation}
    \label{eq:exemplo_p1}
    \int e^x \sin(x) dx = -e^x \cos(x)+ \int e^x \cos(x)dx
  \end{equation}
  Parece que ainda não resolvemos o nosso problema, uma vez que
  $\int e^x \cos(x) dx$ não é uma integral elementar, mas pelo menos
  não é mais complicada. Vamos integrar novamente a integração por
  partes para ver onde chegamos. Dessa vez usaremos $u=e^x$ e
  $dv=\cos(x)dx$. Assim $du=e^x dx$, $v=\sin(x)$, e
  \begin{equation}
    \label{eq:exemplo_p2}
    \int e^x \cos(x) dx= e^x \sin(x)-\int e^x \sin(x) dx
  \end{equation}
  Parece que estamos andando em círculos e que voltamos para o
  mesmo lugar que começamos.
  Porém, se substituirmos a expressão $\int e^x \cos(x)dx$ da
  equação \eqref{eq:exemplo_p2} na
  equação \eqref{eq:exemplo_p1}, temos
  \[
    \int e^x \sin(x) dx
    = -e^x \cos(x)+ e^x \sin(x)-\int e^x \sin(x) dx
  \]
  Assim, temos uma equação para a integral desconhecida e,
  para encontrar o seu valor, basta deixar o termo
  $\int e^x \sin(x) dx$ do lado esquerdo da equação. Assim,
  \[
    2\int e^x \sin(x) dx = -e^x \cos(x)+ e^x \sin(x)
  \]
  Dividindo a equação por 2 e adicionando a constante de integração temos
  \[
    \int e^x \sin(x) dx = \frac{e^x}{2}( -\cos(x) +\sin(x))+ C
  \]
\end{example}
%------------------------------------------------------------------------------%

%------------------------------------------------------------------------------%
\begin{example}
  Encontre \(\int_0^1 \frac{y}{e^{2y}} dy\)
  \solution
  Quando queremos calcular uma integral definida e precisamos usar alguma
  técnica de integração, primeiro calculamos a integral indefinida e depois
  usamos o Teorema Fundamental do Cálculo Parte 2.

  Para encontrar
  \[
    \int \frac{y}{e^{2y}} dy= \int y e^{-2y} dy,
  \]
  usaremos integração por partes e escolhemos
  \begin{align*}
    u  & = y  & dv & = e^{-2y} dy \\
    du & = dy & v  & = -\frac{e^{-2y}}{2}
  \end{align*}
  Portanto,
  \begin{align*}
    \int \frac{y}{e^{2y}} dy
    & =  -\frac{ye^{-2y}}{2}- \int -\frac{e^{-2y}}{2} dy \\
    & = -\frac{ye^{-2y}}{2} + \frac{1}{2} \int e^{-2y} dy \\
    & =  -\frac{ye^{-2y}}{2} - \frac{1}{4} e^{-2y}+ C
  \end{align*}
  Usando o Teorema Fundamental do Cálculo, parte 2
  \begin{align*}
    \int_0^1 \frac{y}{e^{2y}} dy
    & = \left[  -\frac{ye^{-2y}}{2} - \frac{1}{4} e^{-2y} \right] _0^1 \\
    & = \left( -\frac{1}{2}e^{-2}- \frac{1}{4}e^{-2}\right)- \left(0- \frac{1}{4} \right) \\
    & = \frac{1}{4}-\frac{3}{4}e^{-2}
  \end{align*}
\end{example}

\begin{example}
  A figura a seguir mostra a região do primeiro quadrante delimitada pelo eixo
  $y=0$ e pela curva $f(x)=x \sin(x)$, $0 \leq x \leq \pi$.
  \begin{enumerate}[label=\alph*)]
    \item Encontre a área dessa região.
    \item Encontre o volume do sólido de revolução gerado pela
          rotação dessa região em torno do eixo $y$.
    \item Encontre o volume do sólido de revolução gerado pela
          rotação dessa região em torno do da reta $x=\pi$.
  \end{enumerate}
\solution
\begin{enumerate}[label=\alph*)]

\item
  Lembre que a fórmula para calcular a área sob a curva de uma função
  $f(x)$ entre os limites $a$ e $b$ é dada por
  \[
    A = \int_a^b f(x)dx
  \]
  No nosso caso, a função é $f(x) = x \sin(x)$ e os limites são $a=0$ e $b=\pi$.
  Então, a área é
  \[
    A = \int_0^{\pi} x\sin(x)dx
  \]
  Para resolver essa integral definida, precisamos, inicialmente,
  encontrar a integral indefinida $\int x\sin(x)dx$.
  No exemplo \ref{ex:ex1partes} essa integral já foi calculada.
  Assim, usando o Teorema Fundamental do Cálculo-Parte 2
  \[
    A = \left[ -x \cos(x) + \sin(x) \right]_0^{\pi}
    = (-\pi\cos(\pi+ \sin(\pi)))-(-0\cos(0)+ \cos(0))=\pi.
  \]

  \item
  Agora, para encontrar o volume do sólido de revolução gerado pela rotação
  da região delimitada pela curva \(f(x) = x \sin(x)\) e o eixo \(y = 0\)
  em torno do eixo \(y\), usamos o método de cascas cilíndricas, que é
  o mais apropriado. O raio da casca é $x$ e a altura $f(x)$.

  Integrando sobre a região entre \(x = 0\) e \(x = \pi\), obtemos o volume
  \[
    V = \int_{0}^{\pi} 2\pi x \cdot (x \sin(x)) \, dx
      = 2\pi \int_{0}^{\pi} x^2 \sin(x) \, dx
  \]
  Para calcular a integral indefinida, usamos integração por partes.
  Escolhemos \(u = x^2\) e \(dv = \sin(x) \, dx\).
  Então, \(du = 2x \, dx\) e \(v = -\cos(x)\).
  \[
    \int x^2 \sin(x) \, dx = -x^2 \cos(x) - \int -2x \cos(x) \, dx
  \]
  Integramos o segundo termo usando novamente integração por partes.
  Escolhemos \(u = -2x\) e \(dv = \cos(x) \, dx\).
  Então, \(du = -2 \, dx\) e \(v = \sin(x)\).
  \[
    \int -2x \cos(x) \, dx = -2x \sin(x) - \int -2 \sin(x) \, dx
  \]
  \[
    \int -2x \cos(x) \, dx = -2x \sin(x) - 2 \cos(x)+K
  \]
  Agora, voltamos à primeira expressão
  \[
    \int x^2 \sin(x) \, dx = -x^2 \cos(x) - (-2x \sin(x) - 2 \cos(x))+K
  \]

  \[
    \int x^2 \sin(x) \, dx = -x^2 \cos(x) + 2x \sin(x) + 2 \cos(x)+K
  \]
  Agora, usando o Teorema Fundamental do Cálculo-Parte 2
  \begin{align*}
    V & = 2\pi \int_{0}^{\pi} x^2 \sin(x) \, dx                                     \\
      & = 2\pi \left( -x^2 \cos(x) + 2x \sin(x) + 2 \cos(x) \right) \evalat{0}{\pi} \\
      & = 2\pi \Big[ -\pi^2 \cos(\pi) + 2\pi \sin(\pi) + 2 \cos(\pi)                \\
      & \hspace{14mm} - (0^2 \cos(0) + 2(0) \sin(0) + 2 \cos(0)) \Big]
  \end{align*}
  Lembrando que \(\cos(\pi) = -1\) e \(\sin(\pi) = 0\), temos
  \[
    V = 2\pi \left( \pi^2 + 2 \right) = 2\pi^3 - 8\pi
  \]
\end{enumerate}
\end{example}
%------------------------------------------------------------------------------%

%------------------------------------------------------------------------------%
